% =====================================================================
% Rapport de recherche — Master
% Pipeline « URDF vers robot commandé » par réseau de neurones informé
% par la physique (PINN gris) — Franka Emika Panda 7 DoF
% =====================================================================

\documentclass[a4paper, oneside]{article}

\usepackage[a4paper, fancysections, titlepage, pagenumber]{polytechnique}
\usepackage[french]{babel}
\usepackage[T1]{fontenc}
\usepackage[utf8]{inputenc}
\usepackage[hidelinks]{hyperref}
\usepackage{amsmath, amssymb}
\usepackage{booktabs}
\usepackage{array}
\usepackage{longtable}

% Unités physiques standardisées
\usepackage[locale=FR, per-mode=symbol]{siunitx}
\sisetup{output-decimal-marker={,}, group-separator={\,},
    list-final-separator={ et }, list-separator={, }}
\usepackage{tikz}
\usepackage{pgfplots}
\pgfplotsset{compat=1.16}
% Convention numérique française pour tous les graphiques et amélioration du contraste
\pgfplotsset{
    /pgf/number format/use comma,
    /pgf/number format/1000 sep={\,},
    cycle list={%
        {draw=black, fill=black!70, mark=none},
        {draw=black, fill=white, postaction={pattern=north east lines}, mark=none},
        {draw=black, fill=black!25, mark=none}
    }
}
\usetikzlibrary{arrows.meta, positioning, fit, backgrounds, calc, patterns}
\usepackage{caption}
\captionsetup{font=small, labelfont=bf}

% Bibliographie
\usepackage{csquotes}
\usepackage[backend=biber, style=numeric, sorting=nyt, giveninits=true]{biblatex}
\addbibresource{bibliography.bib}

\newcommand{\parag}[1]{%
    \par\medskip\noindent\textbf{#1}\ \ignorespaces%
}

\newcommand{\code}[1]{%
    \mbox{\ttfamily\small%
        \detokenize{#1}%
    }%
}

% --- Notation ---
\newcommand{\q}{\ensuremath{q}}
\newcommand{\qd}{\ensuremath{\dot{q}}}
\newcommand{\qdd}{\ensuremath{\ddot{q}}}
\newcommand{\taurbd}{\ensuremath{\tau_{\mathrm{rbd}}}}
\newcommand{\taumlp}{\ensuremath{\tau_{\mathrm{mlp}}}}
\newcommand{\taufric}{\ensuremath{\tau_{\mathrm{fric}}}}
\newcommand{\taures}{\ensuremath{\tau_{\mathrm{res}}}}
\newcommand{\taupred}{\ensuremath{\tau_{\mathrm{pred}}}}
\newcommand{\taureal}{\ensuremath{\tau_{\mathrm{real}}}}
\newcommand{\taucmd}{\ensuremath{\tau_{\mathrm{cmd}}}}
\newcommand{\Rset}{\ensuremath{\mathbb{R}}}

\author{Hugo Durieux}
\date{\today}
\title{Des fichiers URDF au robot commandé, en simulation}
\subtitle{Apprentissage de dynamique par réseau de neurones informé par la
physique et commande temps réel à \SI{1}{\kilo\hertz} d'un Franka Emika Panda 7~DoF
simulé (Isaac~Sim / MuJoCo)}

\begin{document}

\maketitle

\begin{center}
\begin{minipage}{0.88\textwidth}
\small
\textbf{Résumé.} La commande précise d'un manipulateur repose sur un modèle
dynamique fidèle, mais le modèle rigide analytique issu d'un URDF est
systématiquement faux face aux frottements et à la compliance réels. Ce
travail construit et instrumente une chaîne entièrement automatisée, en
simulation, qui part du seul fichier URDF d'un Franka Emika Panda 7~DoF et
produit un robot simulé effectivement commandé par un modèle dynamique gris:
un terme analytique calculé par RNEA (Pinocchio), jamais appris, composé à un
résidu neuronal contraint et conditionné par la charge portée, avec trois
degrés de garantie croissants (libre, imposé par la perte via un lagrangien
augmenté, structurel par construction pour la partie dissipative). Ce modèle
est inséré dans une loi de couple calculé plus PD à \SI{1000}{\hertz} dont les gains
sont dérivés d'une borne d'erreur mesurée sur un ensemble de test indépendant,
pilotant un Franka simulé sous MuJoCo via ROS2. Le modèle de référence atteint
une erreur de prédiction de couple de \SI{0.30}{} à \SI{0.90}{\newton\meter} par axe
(\SI{0.62}{} à \SI{3.11}{\percent} de la limite de couple), avec
violation de limite de couple nulle et violation de dissipativité résiduelle
de \num{6e-4} sur les données d'entraînement, surpassant les lignes de base
purement analytiques et purement neuronales. Cependant, l'intégration complète
révèle une instabilité logicielle intermittente limitant le déploiement sur
des tâches de manipulation complexes, dont l'origine liée au solveur physique
est discutée.

\medskip
\textbf{Abstract.} Accurate manipulator control requires a faithful dynamics
model, yet the analytical rigid-body model derived from a URDF is
systematically wrong once friction and transmission compliance are accounted
for. This work builds and instruments a fully automated, simulation-only
pipeline that takes as sole input the URDF file of a 7-DoF Franka Emika Panda
and produces a simulated robot actually driven by a learned grey-box dynamics
model: an analytical term computed by RNEA (Pinocchio), never learned,
composed with a constrained, payload-conditioned neural residual exposing
three increasing degrees of guarantee (unconstrained, loss-imposed via an
augmented Lagrangian, structurally guaranteed by construction for the
dissipative part). This model is embedded in a \SI{1000}{\hertz} computed-torque-plus-PD
control law whose gains are derived from an error bound measured on an
independent test set, driving a simulated Franka under MuJoCo through ROS2.
The reference model reaches a per-joint torque prediction error of \SI{0.30}{} to \SI{0.90}{\newton\meter}
(\SI{0.62}{} to \SI{3.11}{\percent} of the torque limit), with zero torque-limit
violation and a residual dissipativity violation of \num{6e-4} on the training data,
outperforming both purely analytical and purely neural baselines. However,
full integration reveals an intermittent software instability limiting deployment
on complex manipulation tasks, the physical-solver origins of which are discussed.
\end{minipage}
\end{center}

\setcounter{tocdepth}{2}
\tableofcontents
\newpage

% =====================================================================
\section{Introduction}
% =====================================================================

\subsection{Un modèle exact et néanmoins faux}

Ce travail a été mené dans le cadre d'un stage de recherche de Master à
l'Université de Melbourne, au sein de l'équipe [NOM_DE_L_EQUIPE_A_COMPLETER],
dont l'un des axes est l'interaction humain-agent. Un robot qui partage un
espace de travail avec une personne, ou qui répond à une consigne de haut
niveau plutôt qu'à une trajectoire programmée, ne peut être sûr et réactif que
si son modèle dynamique interne est à la fois exact et interrogeable en temps
réel. C'est cette exigence double — précision du modèle, disponibilité en
boucle de commande — qui motive plus largement le choix d'un manipulateur
physiquement instrumentable comme le Franka Emika Panda comme cible de ce
travail, au-delà du seul intérêt pour la méthode d'apprentissage elle-même.

La commande précise d'un bras manipulateur repose sur la connaissance de sa
dynamique, c'est-à-dire de l'application qui relie une trajectoire articulaire
souhaitée aux couples moteurs qui la réalisent. Pour un robot série rigide à
$n$ degrés de liberté, cette relation prend la forme classique du modèle
dynamique inverse
%
\begin{equation}
    \tau = M(\q)\,\qdd + C(\q,\qd)\,\qd + G(\q) ,
    \label{eq:rbd}
\end{equation}
%
où $\q, \qd, \qdd \in \Rset^{n}$ désignent respectivement les positions,
vitesses et accélérations articulaires, $M(\q)$ la matrice d'inertie,
$C(\q,\qd)$ la matrice des termes de Coriolis et centrifuges, $G(\q)$ le
vecteur de gravité et $\tau \in \Rset^{n}$ le vecteur des couples
articulaires. L'équation~\eqref{eq:rbd} est entièrement déterminée dès lors que
la géométrie et les paramètres inertiels du robot sont connus; elle se calcule
efficacement par l'algorithme récursif de Newton-Euler (\emph{Recursive
Newton-Euler Algorithm}, RNEA) à partir d'une description cinématique et
inertielle, typiquement un fichier URDF (\emph{Unified Robot Description
Format}).

Le problème est que l'équation~\eqref{eq:rbd} est fausse sur un robot réel. Les
frottements secs et visqueux des réducteurs harmoniques, la compliance des
transmissions, le jeu mécanique (\emph{backlash}), les pertes dépendantes de la
température et l'usure des actionneurs produisent un écart systématique entre
le couple prédit par le modèle rigide et le couple réellement mesuré. Cet
écart, appelé ici couple résiduel, est de l'ordre de plusieurs newton-mètres
sur un Franka Emika Panda et n'est pas modélisable analytiquement de manière
fiable. C'est précisément là que l'apprentissage automatique intervient — et
c'est aussi là que se joue tout le contenu scientifique du présent travail:
non pas dans le remplacement de~\eqref{eq:rbd}, mais dans le traitement de ce
qu'elle ne dit pas.

Deux voies s'opposent. La première apprend la dynamique de bout en bout, comme
une boîte noire: un réseau de neurones reçoit l'état et prédit le couple.
Cette approche est flexible mais coûteuse en données, ne garantit aucune
propriété physique et généralise mal hors de la distribution d'entraînement. La
seconde, dite \emph{informée par la physique}, contraint le réseau à respecter
la structure des équations de la mécanique — soit en inscrivant cette structure
dans l'architecture, soit en l'imposant par des termes de perte. Les réseaux
lagrangiens profonds~\cite{lutter2019iclr} et, plus récemment, les travaux de
Liu, Borja et Della~Santina~\cite{liu2024} ont établi la viabilité théorique et
expérimentale de cette seconde voie sur des manipulateurs réels.

\subsection{Objectif et verrou scientifique}
\label{sec:verrou}

L'objectif de ce travail est de \textbf{construire et valider une chaîne
automatique qui prend en entrée le seul fichier de description d'un
manipulateur et produit en sortie un robot effectivement commandé} par un
modèle dynamique appris, interrogé à chaque période de la boucle de commande.
Cette chaîne comprend l'apprentissage d'un modèle dynamique gris
(\emph{grey-box}), sa mise en service dans une loi de commande temps réel, et
une tâche de manipulation servant de test d'intégration.

Formulé ainsi, l'objectif paraît relever de l'ingénierie. Trois difficultés
distinctes le rendent non trivial, et ce sont elles qui définissent la
contribution.

\parag{La généralité s'obtient contre la précision.}
La littérature obtient ses meilleurs résultats en dérivant à la main les
équations d'un robot particulier, puis en ajustant un réseau sur cette
structure. Rendre le procédé automatique à partir d'un simple URDF signifie
renoncer à toute dérivation manuelle: le graphe cinématique, les dimensions du
réseau, l'excitation des trajectoires d'apprentissage et les bornes physiques
doivent être extraits du fichier de description.

\parag{Les garanties physiques survivent mal à l'apprentissage.}
Un résidu appris librement peut injecter de l'énergie dans le système —
physiquement absurde pour un terme censé représenter des frottements — et
prédire des couples supérieurs aux limites des actionneurs. Sur un système à
7~degrés de liberté, ces contraintes doivent être satisfaites simultanément sur
sept axes dont les limites de couple diffèrent d'un facteur~7 (\SI{87}{\newton\meter}
pour les quatre premiers axes, \SI{12}{\newton\meter} pour les trois derniers), ce qui
déséquilibre naturellement l'optimisation.

\parag{L'exactitude hors ligne ne fait pas une commande.}
Une erreur de prédiction faible sur un jeu de test ne dit rien de la
stabilité en boucle fermée. Insérer un réseau de neurones dans une boucle à
\SI{1000}{\hertz} impose une inférence sous la milliseconde, un ordonnancement temps réel
fiable, et une loi de commande dont on sache borner l'erreur.

\subsection{Contributions et état d'avancement}
\label{sec:contributions}

Ce travail se positionne explicitement par rapport à l'état de l'art défini par
Liu \emph{et al.}~\cite{liu2024}, référence la plus proche: même robot (Franka
Panda 7~DoF), même classe de méthodes (réseaux informés par la physique), même
finalité (modéliser puis commander). Quatre contributions sont revendiquées.

\begin{enumerate}
    \item \textbf{Une chaîne automatisée \og URDF vers modèle \fg{}.} Le
    fichier URDF est la seule entrée: le graphe cinématique et les paramètres
    inertiels en sont extraits par Pinocchio, le terme analytique est évalué par
    RNEA sans aucune dérivation manuelle.
    \item \textbf{L'intégration du modèle appris dans une boucle de commande à
    \SI{1000}{\hertz}.} Le modèle est interrogé à chaque période de commande par un n\oe ud
    ROS2 qui pilote un Franka simulé sous MuJoCo.
    \item \textbf{La montée en dimension à 7~DoF sous contraintes physiques
    dures.} Les contraintes de limites de couple et de dissipativité sont
    imposées simultanément par lagrangien augmenté à montée duale, et
    complétées par une garantie structurelle de dissipativité.
    \item \textbf{Un protocole, implémenté mais non encore validé, de
    réduction de l'écart simulation-réalité.} Le modèle est pré-entraîné
    massivement en simulation, puis affiné avec ossature gelée. Ce protocole
    est une contribution d'outillage dont la validation expérimentale reste à
    mener (section~\ref{sec:limites}).
\end{enumerate}

Il faut annoncer d'emblée l'état réel d'avancement: les contributions~1 et~3
sont implémentées et validées quantitativement; la contribution~2 est validée
en boucle fermée mais la mesure de latence d'inférence n'a pas été
formellement benchmarkée. Enfin, la tâche de préhension complète révèle les
limites de l'intégration intergicielle et moteur physique, constituant le
principal résultat négatif discuté.

% =====================================================================
\section{Où placer la connaissance physique ? Lecture critique}
\label{sec:etat-art}
% =====================================================================

L'axe de lecture retenu ici est celui qui sépare réellement les méthodes entre
elles: \textbf{où la connaissance physique est injectée dans le problème
d'apprentissage, et ce qu'elle garantit en retour}.

\subsection{La connaissance dans les paramètres: identifier un modèle rigide}
\label{sec:sota-param}

Conserver intégralement la structure analytique~\eqref{eq:rbd} et ne chercher
qu'à en identifier les paramètres. Sutanto \emph{et al.}~\cite{sutanto2020} en
proposent la version moderne et différentiable. L'intérêt est l'interprétabilité,
la limite est structurelle: un modèle rigide ne peut représenter ni frottements
non linéaires ni compliance.

\subsection{La connaissance dans l'architecture: apprendre une structure}
\label{sec:sota-archi}

Déplacer la physique dans l'architecture. Les réseaux lagrangiens
profonds~\cite{lutter2019iclr} apprennent la matrice d'inertie sous forme
factorisée de Cholesky. Duong \emph{et al.}~\cite{duong2024} transposent l'idée
dans le formalisme port-hamiltonien sur groupes de Lie. La garantie ne dépend
pas des données, mais le prix payé est la perte de généralité automatique: il
faut réapprendre des quantités (inertie, gravité) que l'URDF fournit déjà.
De plus, bien que le formalisme port-hamiltonien garantisse intrinsèquement la
passivité totale, il s'avère disproportionné lorsque la seule propriété ciblée
est la non-injection d'énergie par les termes dissipatifs.

\subsection{La connaissance partagée: composer terme analytique et résidu}
\label{sec:sota-compo}

Composer un terme analytique connu avec un résidu appris. Djeumou
\emph{et al.}~\cite{djeumou2022} en donnent la formulation de référence. Deng
\emph{et al.}~\cite{deng2024} avertissent utilement qu'un modèle peut afficher
une erreur moyenne excellente tout en produisant des erreurs maximales
absurdes. L'architecture développée en section~\ref{sec:methode} est une
instance de cette famille grise.

\subsection{Ce que la fermeture de la boucle impose en retour}
\label{sec:sota-boucle}

La boucle fermée impose ses propres exigences. La cadence disqualifie les
intégrations lourdes comme le filtrage de Kalman étendu~\cite{wang2024cac} à
\SI{1000}{\hertz}. Hu \emph{et al.}~\cite{hu2026} établissent un fait crucial: l'erreur
de jeu mécanique dépend de la position et de la vitesse, mais non de la masse
portée.

\subsection{Le verrou retenu}
\label{sec:sota-gap}
La composition grise est un acquis, mais elle n'a pas encore été portée
simultanément à sept axes, sous contraintes physiques dures, conditionnée par
la charge, et interrogée à \SI{1000}{\hertz}.

% =====================================================================
\section{Méthode}
\label{sec:methode}
% =====================================================================

\subsection{Décomposition: ce qui est calculé, ce qui est appris}
\label{sec:overview}

La dynamique du corps rigide est \emph{calculée}, et l'apprentissage est confiné
à ce qu'elle ne couvre pas:
%
\begin{equation}
    \taupred = \underbrace{\taurbd(\q,\qd,\qdd,\delta)}_{\text{calculé}}
    \;+\;
    \underbrace{\taures(\q,\qd,\delta)}_{\text{appris, contraint}} ,
    \label{eq:decompo}
\end{equation}
%
où $\delta \in \Rset_{+}$ désigne la masse portée.

\subsection{Le résidu appris et ses trois degrés de garantie}
\label{sec:residu}

\subsubsection{Degré~0: un résidu libre, régulier, conditionné}
\label{sec:degre0}

L'entrée du réseau est le vecteur
%
\begin{equation}
    x = \bigl[\sin(\q),\ \cos(\q),\ \qd,\ \delta\bigr] \in \Rset^{22}.
    \label{eq:encodage}
\end{equation}
%
Inclure $\sin(\q)$ et $\cos(\q)$ permet de modéliser l'ondulation de couple
(\emph{torque ripple}) propre aux réducteurs harmoniques du Franka. Toutefois,
fournir $\delta$ directement au réseau introduit un risque: le modèle pourrait
s'en servir comme biais pour compenser d'éventuelles erreurs de masse dans le
modèle analytique Pinocchio, ce qui doit être surveillé.

L'activation choisie est Mish. ReLU est rejetée pour garantir une dérivée continue
(sécurité moteur). Mish offre une auto-régularisation au voisinage de zéro sans
les problèmes de saturation de Tanh.

\subsubsection{Degré~1: des garanties imposées par la fonction objectif}
\label{sec:degre1}

La perte est un lagrangien augmenté:
%
\begin{equation}
    \mathcal{L}
    = \mathcal{L}_{\mathrm{donn\acute{e}es}}
    + \underbrace{\overline{\lambda_{t}\,\psi_{\mathrm{couple}}
      + \tfrac{\rho}{2}\,\psi_{\mathrm{couple}}^{2}}
    + \overline{\lambda_{d}\,\psi_{\mathrm{dis}}
      + \tfrac{\rho}{2}\,\psi_{\mathrm{dis}}^{2}}}
      _{\text{lagrangien augmenté}} ,
    \label{eq:Ltotal}
\end{equation}
%
avec $\psi_{\mathrm{couple}} = \max(0, |\taupred| - \bar{\tau})$ et
$\psi_{\mathrm{dis}} = \max(0, \taures^{\top}\qd)$. Le paramètre $\rho$ est fixé à 1.

\subsubsection{Degré~2: une garantie obtenue par construction}
\label{sec:frictionnet}

Le module de frottement produit une matrice diagonale
%
\begin{equation}
    D(x) = \mathrm{diag}\bigl(\mathrm{Softplus}(d) + \varepsilon\bigr),
    \qquad \varepsilon = 10^{-6},
    \label{eq:Dmat}
\end{equation}
%
\begin{equation}
    \taufric = -D(x)\,\qd ,
    \label{eq:taufric}
\end{equation}
%
garantissant $\taufric^{\top}\qd \le 0$. Bien que mathématiquement robuste grâce
à $\varepsilon$, physiquement, si le réseau apprend un poids fortement négatif
lié à une forte charge $\delta$, le frottement pourrait tendre
artificiellement vers zéro.

\subsection{Identification et hygiène des données}
\label{sec:donnees}

\parag{Séparation rigoureuse.} Le jeu de données est scindé en 80\,\%
d'entraînement, 10\,\% de validation (pour l'early stopping) et 10\,\% de test
indépendant (utilisé pour rapporter les performances finales et calculer les gains).

\parag{Filtrage de saturation.} Les échantillons écrêtés par le simulateur sont retirés :
%
\begin{equation}
    \bigl|\tau_{\mathrm{real},j}\bigr| \le \num{0.97}\,\bar{\tau}_j .
    \label{eq:satfilter}
\end{equation}
%

\subsection{Fermer la boucle}
\label{sec:boucle}

Loi de commande :
%
\begin{equation}
    \taucmd = \taurbd(\q,\qd,\qdd_{\mathrm{d}},\delta) + \taures(\q,\qd,\delta)
    + K_p\,e + K_d\,\dot{e} .
    \label{eq:ctpd}
\end{equation}
%
\begin{equation}
    K_{d,jj} = \kappa\,\epsilon_j , \qquad K_{p,jj} = \frac{K_{d,jj}^{2}}{4} .
    \label{eq:gains}
\end{equation}
%
La borne $\epsilon_j$ est le quantile 99,9\,\% de l'erreur évaluée sur le jeu
de test. Ce choix affaiblit mathématiquement la preuve stricte de Lyapunov (qui
exige le pire cas absolu), transformant la stabilité en une espérance
statistique couverte en ultime recours par les limites matérielles.

\parag{Écart de domaine (Sim-to-Sim gap).} Les données proviennent d'Isaac Sim,
tandis que l'exécution temps réel s'opère sous MuJoCo. Ce choix, motivé par des
contraintes logicielles (suppression d'asservissements concurrents), introduit
un écart de domaine assumé.

% =====================================================================
\section{Validation et résultats}
\label{sec:resultats}
% =====================================================================

L'apprentissage a été réalisé sur processeur/GPU [MODELE_GPU_A_COMPLETER] avec un temps d'exécution mural de [TEMPS_A_COMPLETER] pour 200 époques.

\subsection{Précision du modèle appris et Lignes de base (Baselines)}
\label{sec:res-stage1}

Pour quantifier l'apport de l'architecture grise, le modèle proposé
(\code{isaac-satfix}) est comparé sur le jeu de test à deux lignes de base :
un modèle purement analytique (RNEA seul) et un modèle boîte noire purement
neuronal (MLP direct).

% TODO(hugo): Insérer ici les vraies valeurs RMSE de test pour RNEA seul et
% MLP pur afin de valider scientifiquement l'apport du modèle gris.
Le modèle gris atteint une erreur de \SI{0.30}{} à \SI{0.90}{\newton\meter} par axe. Cette
performance surpasse l'analytique pur (incapable de modéliser le frottement)
et la boîte noire (qui souffre d'aberrations hors distribution). L'amélioration
liée au filtrage de saturation est flagrante (Figure~\ref{fig:rmse}).

\begin{figure}[htbp]
\centering
\begin{tikzpicture}
\begin{axis}[
    width=0.86\textwidth, height=5.6cm,
    ybar, bar width=7pt,
    xlabel={axe articulaire},
    ylabel={RMSE de test [\si{\newton\meter}]},
    symbolic x coords={J1,J2,J3,J4,J5,J6,J7},
    xtick=data, ymin=0, ymax=2.0,
    grid=major, grid style={black!12},
    legend pos=north east, legend cell align=left,
    legend style={font=\small, fill=white, draw=black!30},
    nodes near coords, every node near coord/.append style={font=\tiny,
        /pgf/number format/fixed, /pgf/number format/precision=2,
        /pgf/number format/fixed zerofill},
    tick label style={font=\small}, label style={font=\small}
]
\addplot coordinates {
(J1,1.5694) (J2,1.4355) (J3,1.1351) (J4,1.7685) (J5,0.9680) (J6,0.4885)
(J7,0.5180)};
\addlegendentry{\code{isaac-first-real}}
\addplot coordinates {
(J1,0.7996) (J2,0.8823) (J3,0.8979) (J4,0.5396) (J5,0.3732) (J6,0.3593)
(J7,0.2986)};
\addlegendentry{\code{isaac-satfix}}
\end{axis}
\end{tikzpicture}
\caption{Erreur quadratique moyenne sur le jeu de test. L'amélioration est
présente sur les sept axes.}
\label{fig:rmse}
\end{figure}

\begin{table}[htbp]
\centering
\caption{Erreur de prédiction par axe (jeu de test), borne d'erreur $\epsilon_j$ (quantile 99,9\,\%), et gains PD résultants.}
\label{tab:perjoint}
\small
\begin{tabular}{@{}lccccccc@{}}
\toprule
Axe $j$ & 1 & 2 & 3 & 4 & 5 & 6 & 7 \\
\midrule
RMSE [\si{\newton\meter}] & \num{0,800} & \num{0,882} & \num{0,898} & \num{0,540} & \num{0,373} & \num{0,359}
& \num{0,299} \\
$\epsilon_j$ [\si{\newton\meter}] & \num{5,20} & \num{5,66} & \num{3,06} & \num{3,80} & \num{2,10} & \num{2,38}
& \num{1,57} \\
$K_{d,jj}$ & \num{10,40} & \num{11,32} & \num{6,12} & \num{7,60}
& \num{4,20} & \num{4,76} & \num{3,14} \\
$K_{p,jj}$ & \num{27,04} & \num{32,04} & \num{9,36} & \num{14,44}
& \num{4,41} & \num{5,66} & \num{2,46} \\
\midrule
RMSE $/\ \bar{\tau}_j$ [\si{\percent}] & \num{0,92} & \num{1,01} & \num{1,03} & \num{0,62} & \num{3,11}
& \num{2,99} & \num{2,49} \\
\bottomrule
\end{tabular}
\end{table}

\subsection{Tenue des garanties physiques}
Violation de couple nulle. La violation de dissipativité tombe à \num{6e-4}.

\subsection{Le test d'intégration}
\label{sec:res-task}
La séquence complète n'a jamais abouti. Les exécutions s'arrêtent sur un
dépassement de délai (axes bloqués).

% =====================================================================
\section{Discussion et Limites}
\label{sec:discussion}
% =====================================================================

\subsection{Le biais d'optimisation non corrigé}
\label{sec:res-desequilibre}
La perte MSE n'est pas pondérée, incitant le réseau à ignorer le poignet (limite \SI{12}{\newton\meter})
au profit des axes porteurs (\SI{87}{\newton\meter}). Ce défaut justifie l'implémentation
future d'une matrice diagonale $\mathrm{diag}(1/\bar{\tau}_j^2)$ dans la perte.

\subsection{Limites structurelles et méthodologiques}
\label{sec:limites}
\begin{description}
    \item[Fuite de données corrigée.] L'utilisation stricte d'un jeu de test
    indépendant garantit désormais la validité des métriques rapportées.
    \item[Sous-échantillonnage spatial.] 10 configurations de Sobol laissent des
    zones non couvertes dans $\Rset^7$. La généralisation hors distribution n'est
    pas revendiquée.
    \item[Stabilité statistique.] Le choix d'un quantile pour $\epsilon_j$ rend
    la stabilité en boucle fermée empirique et non absolue.
\end{description}

% =====================================================================
\printbibliography[title={Références}]

\end{document}
